\documentclass[12pt]{article}

\usepackage[utf8]{inputenc}
\usepackage{latexsym,amsfonts,amssymb,amsthm,amsmath}
\usepackage{tikz}
\usepackage{stanli}
\usetikzlibrary{shapes,calc,through,3d,patterns}
\usetikzlibrary{arrows.meta,decorations.pathreplacing,angles,quotes}
\usepackage{multicol}
\usepackage{geometry}
\usepackage{mathtools}
\usepackage{siunitx}
\usepackage{booktabs}

\setlength{\parindent}{0in}
\setlength{\oddsidemargin}{-0.25in}
\setlength{\textwidth}{7in}
\setlength{\textheight}{9.3in}
\setlength{\topmargin}{-1in}
\setlength{\headheight}{18pt}

\newcommand\blfootnote[1]{%
    \begingroup
    \renewcommand\thefootnote{}\footnote{#1}%
    \addtocounter{footnote}{-1}%
    \endgroup
}

\definecolor{darkred}{rgb}{0.6,0,0}

\newcommand{\bolddelta}{\boldsymbol\delta}
\newcommand{\boldxi}{\boldsymbol\xi}

\begin{document}

\noindent ME473 - Dynamic finite element analysis of structures\hfill EPFL\\
Stefano Burzio \hfill 2025\\
\vspace{0.3cm}

\hspace*{\fill}\textbf{\Large Lecture notes week 2}\hspace*{\fill}

\hrulefill


\subsubsection*{Semi-discrete formulation of elastodynamics.} In linear elastodynamics, the displacement field \( \mathbf  \mathbf u(\mathbf x, t) \) satisfies the strong form:
\[
    \nabla^T \mathbf C \nabla\mathbf  u + \mathbf f = \rho \ddot{\mathbf u},
\]
where \(\mathbf  C \) is the elasticity tensor, \( \mathbf f \) the body force, and \( \rho \) the density.

We obtain the weak form by multiplying by a virtual displacement \( \bolddelta \mathbf u \) and integrating over the domain:
\[
    \int_\Omega (\nabla \bolddelta u)^T \mathbf C \nabla \mathbf u \, d\Omega + \int_\Omega \rho \bolddelta \mathbf u^T \ddot{\mathbf u} \, d\Omega = \int_\Gamma \bolddelta \mathbf u^T \hat{\mathbf f} \, d\Gamma + \int_\Omega \bolddelta \mathbf u^T \mathbf f \, d\Omega.
\]

Using the Galerkin method, we approximate:
\begin{itemize}
    \item \(\mathbf u(\mathbf x,t) \approx\mathbf u^h(\mathbf x,t) =\mathbf H(\mathbf x) \mathbf q(t) \),
    \item \( \bolddelta\mathbf u(\mathbf x) \approx \bolddelta\mathbf u^h(\mathbf x) = \mathbf H(\mathbf x) \bolddelta \mathbf q \),
\end{itemize}
where \( \mathbf H(\mathbf x) \) contains the shape functions and \( \mathbf q(t) \) the nodal degrees of freedom.

This yields the semi-discrete weak form:
\[
    \mathbf M \ddot{\mathbf q}(t) + \mathbf K \mathbf q(t) = \mathbf r(t),
\]
with:
\begin{align*}
    \mathbf M    & = \int_\Omega \rho \mathbf H^T \mathbf H \, d\Omega,                                                 \\
    \mathbf K    & = \int_\Omega \mathbf B^T\mathbf C\mathbf B \, d\Omega,                                              \\
    \mathbf r(t) & = \int_\Gamma \mathbf H^T \hat{\mathbf f} \, d\Gamma + \int_\Omega \mathbf H^T \mathbf f \, d\Omega.
\end{align*}
Here, \( \mathbf B = \nabla\mathbf H \) is the strain-displacement matrix.

\subsubsection*{Initial conditions}
Given \( \mathbf u(\mathbf x,0) = \mathbf u_0(\mathbf x) \), \( \dot{\mathbf u}(\mathbf x,0) = \mathbf v_0(x) \), the initial values for \( \mathbf q(t) \) and \( \dot{\mathbf q}(t) \) are:
\[
    \mathbf q(0) =\mathbf u_0(\mathbf x_j), \quad \dot{\mathbf q}(0) = \mathbf v_0(\mathbf x_j).
\]

\subsubsection*{Boundary conditions}
On \( \Gamma_u \), we enforce:
\[
    \mathbf q_j(t) = \hat{\mathbf u}(\mathbf x_j,t), \quad \bolddelta\mathbf  q_j = 0.
\]

\subsubsection*{Advantages}
\begin{itemize}
    \item Systematic approximation method.
    \item Reuses the same shape functions for test and trial functions.
    \item Converges efficiently with suitable shape functions.
\end{itemize}

\subsubsection*{Drawbacks}
\begin{itemize}
    \item Heavily dependent on the chosen basis.
    \item No direct physical interpretation of \( q(t) \).
    \item Boundary and initial conditions may be cumbersome in practice.
\end{itemize}

\end{document}